\section{Introduction}\label{sec:introduction}

Our goal was to measure how much a land-use classifier trained on European Sentinel-2 images
loses when it is used in Aguascalientes, a semi-arid state in central Mexico, and how much of
that loss a small set of local labels can win back. We compared two convolutional neural
networks (CNNs), ResNet-50 and EfficientNet-B0, with a Vision Transformer, ViT-S/16. All three
were trained on EuroSAT and then tested on patches of Aguascalientes that we labeled by hand.

Land use matters in Aguascalientes because of water. The state water plan names the
overexploitation of groundwater, the main source of supply, as the central problem of the
state~\cite{inagua2022planhidrico}. The same plan states that the demand for water comes mainly
from agricultural and urban public use. Both uses can be seen from a satellite: farmland and the
city are land-use classes that a Sentinel-2 image can show. Monitoring land use is therefore one
way to follow the pressure on water, although it does not measure the water itself.

Monitoring land use in Aguascalientes faced a practical problem. There were no labeled Sentinel-2 patches of Aguascalientes,
and the land-use models available had been trained on European landscapes such as those of
EuroSAT. The long-term goal of the project is a monthly land-use monitor of the state, which is
the second stage of this work. That monitor needs a model that works in Aguascalientes, so this
paper chooses one and measures how well it transfers.

The difficulty has a name: geographic domain shift. A model learns how each class looks in the
region where it was trained, and a new region can show the same class with other colors and
textures. A model that learned Forest as a dark, closed canopy and Pasture as green grass has
little reason to recognize their versions in semi-arid Aguascalientes, imaged in April, at the
peak of the dry season. The size of this loss can only be measured
fairly on a test set that is spatially separated from the labels used to adapt the model.

We organized the study around five research questions. Which of ResNet-50, EfficientNet-B0 and
ViT-S/16 performs best on EuroSAT, and at what computational cost (RQ-a)? How much does each
model degrade on Aguascalientes without any local training, that is, zero-shot (RQ-b)? How much
of that degradation does light fine-tuning recover (RQ-c)? Does 13-band input outperform RGB
(RQ-d)? And do self-supervised SSL4EO-S12 weights outperform ImageNet weights (RQ-e)?

The short answer is that the domain shift outweighed the differences between the models. On
EuroSAT, the eight runs reached a macro-F1, the average of the per-class F1 scores, between
\EuroSATMin{} and \EuroSATMax. In Aguascalientes, without local training, they lost between
\DropMin{} and \DropMax{} in macro-F1, and fine-tuning the last block recovered at most
\RecoveryMax{} of that drop (Eq.~\ref{eq:recovery}). The highest observed macro-F1 after fine-tuning,
\AgsLastBlockBest, came from ResNet-50 with 13 bands and SSL4EO-S12 weights, the model we
selected.

Our contributions are the following:
\begin{itemize}
  \item A comparison of ResNet-50, EfficientNet-B0 and ViT-S/16 under one shared training
  protocol, including RGB versus 13-band input and ImageNet versus SSL4EO-S12 weights.
  \item A measurement of the drop from EuroSAT to Aguascalientes and of its partial recovery
  with two light fine-tuning modes, on \num{1102} hand-labeled patches (561 for fine-tuning and
  541 for testing) split by blocks of 8$\times$8 cells (5.12~km).
  \item Open code, labels and checkpoints, so that the experiments can be repeated and the
  models reused.
\end{itemize}

The rest of the paper is organized as follows. Section~\ref{sec:related} reviews related work.
Section~\ref{sec:data} describes the two datasets and the spatial split, and
Section~\ref{sec:methodology} the models, the shared protocol and the metrics.
Section~\ref{sec:results} reports the results, Section~\ref{sec:discussion} discusses their
likely causes and limits, and Section~\ref{sec:conclusions} concludes.
